\documentclass[11pt]{article}
\usepackage[margin=1in]{geometry}
\usepackage{amsmath,amssymb,amsthm}
\newtheorem{theorem}{Theorem}
\newtheorem{lemma}{Lemma}
\newtheorem{corollary}{Corollary}
\newcommand{\bits}{\{0,1\}}
\begin{document}
\title{Laboratory V97: Comparison-Free Peeling Kernels}
\author{NC0\_k-Avoid Laboratory}
\date{2026-08-10}
\maketitle

Let $C:\bits^N\to\bits^{N+1}$ be an explicit locality-three circuit.  Replace
every listed support by its essential support.  In the bipartite incidence graph,
a connected component $K$ has $n_K$ inputs and $m_K$ outputs.  Since
$\sum_K(m_K-n_K)=1$, some component has positive surplus.

\begin{lemma}[Leaf-pair reduction]
If an input $x$ occurs in exactly one active output $g$, delete both $x$ and
$g$.  Every missing word of the reduced output map lifts to a missing word of
the original map after assigning an arbitrary target bit to $g$.
\end{lemma}
\begin{proof}
No surviving output depends on $x$.  Hence any original assignment realizing
the lifted target would realize the reduced missing word on the surviving
outputs, a contradiction.  One input and one output are deleted, so surplus is
preserved.
\end{proof}

\begin{lemma}[Unary forcing]
If an active output $g$ has essential support $\{x\}$, choose target $g=0$.
There is a unique $a\in\bits$ with $g(a)=0$.  Delete $g$ and $x$, substitute
$x=a$ into all surviving gates, and normalize again.  Every missing word of the
restricted circuit lifts to a missing word of the pre-restriction circuit.
\end{lemma}
\begin{proof}
Any assignment realizing target $g=0$ must have $x=a$, so its other outputs are
exactly outputs of the restricted circuit.  Surplus is again preserved.
\end{proof}

Also delete unused inputs, and if a constant output appears choose the opposite
target bit and terminate.  Apply these operations in a fixed deterministic
order.  For a positive-surplus component $K$, let $\lambda(K)$ be the number of
inputs remaining in the nonterminated kernel, or zero after constant
termination, and put
\[
 \lambda(C)=\min_{m_K>n_K}\lambda(K).
\]

\begin{theorem}[Peeling-kernel avoider]
There is a deterministic comparison-free algorithm producing a word outside
$\mathrm{Range}(C)$ in time
\[
 O\!\left(2^{\lambda(C)}\operatorname{poly}(N)\right).
\]
In particular the algorithm is polynomial when $\lambda(C)=O(\log N)$.
\end{theorem}
\begin{proof}
Every nonterminating reduction preserves or increases positive surplus, so a
final kernel with $\lambda$ inputs has more than $\lambda$ outputs.  Enumerate
its $2^\lambda$ input assignments and hence its local range.  Among any
$2^\lambda+1$ distinct kernel-output words one is absent.  Reverse all recorded
leaf and unary reductions; the two lemmas preserve absence at each step.  Fill
outputs outside the selected component arbitrarily.
\end{proof}

Let $\rho(C)$ be the V96 parameter: the minimum $n_K$ over positive-surplus
components.

\begin{corollary}
For every circuit, $\lambda(C)\leq\rho(C)$.
\end{corollary}
\begin{proof}
The reducer only deletes inputs from each candidate component.
\end{proof}

\begin{theorem}[Strict separation]
For every $N\ge8$ there is a connected, exact-stretch, nonmonotone circuit with
genuinely ternary gates satisfying
\[
 \rho(C_N)=N,
 \qquad
 \lambda(C_N)=\lceil\log_2 N\rceil.
\]
\end{theorem}
\begin{proof}
Let $h=\max\{3,\lceil\log_2N\rceil\}$.  Put $h+1$ parity-of-three outputs on
cyclic triples of $h$ core inputs.  For every remaining input $z$, add one
parity-of-three output on $(x_0,x_1,z)$.  The full incidence graph is connected,
so $\rho=N$.  Every attachment $z$ has input degree one and is peeled with its
unique output, leaving exactly the $h$-input, $h+1$-output parity core.  That
core has no applicable reduction, hence $\lambda=h$.
\end{proof}

The result is parameterized and does not improve the known worst-case
$O(N2^{N/2})$ locality-three bound, does not solve unrestricted range avoidance,
and does not resolve P versus NP.
\end{document}
